\section{Step 4 Finish exactly sixteen scalar rows}\label{h263-sad-s04}
Advance both sample cursors by one after each pair, and repeat for sixteen pairs.
When a row finishes, advance the retained row starts by their respective strides.
The row starts and moving cursors have different jobs: adding the stride to a
cursor already sixteen bytes past the start would drift sixteen extra samples
on every row.

Finish sixteen rows and retain the complete sum. For the worked first row, eight
differences are 144 and eight are 132, giving 2208. Later rows overlap the bright
square differently, so multiplying 2208 by sixteen is not the block result.
The complete scalar sum is 7872. A row's neighbors and padding do not contribute.

\paragraph{Observe} The successful trace has 256 sample loads from each window.
Their addresses fit the sixteen active bytes in each row. The final sum is 7872.
\Needspace{8\baselineskip}
\paragraph{Check} Which pointer should receive the row stride?
\begin{itemize}
  \item[A.] The pointer to the start of the current row.
  \item[B.] The pointer already sixteen samples beyond that start, without adjustment.
  \item[C.] The address of the returned sum.
\end{itemize}
